\section{Description des responsabilités}

\subsection{Détail des attributions par acteur}

Conformément aux éléments communiqués lors de la réunion de cadrage, les responsabilités métier et les opérations associées à chaque acteur s'établissent comme suit :

\subsubsection{1. L'Ordonnateur}
L'ordonnateur assure le pilotage administratif de la dépense de son initialisation jusqu'à l'ordre de payer. Ses attributions portent sur :
\begin{itemize}
  \item l'initiative de la dépense par l'expression du besoin ;
  \item la formalisation de l'engagement juridique (préparation du projet d'acte juridique, soumission au contrôleur financier, puis signature de l'acte après obtention des contrôles requis) ;
  \item la prise en compte de l'exécution physique à travers la réception des informations relatives aux livraisons transmises par le comptable matière ;
  \item la vérification des droits du créancier et l'arrêté du montant de la dépense (enregistrement et validation des factures, réalisation de la liquidation) ;
  \item l'émission de l'ordre de paiement (génération du projet d'ordonnancement de paiement, signature et transmission au contrôleur financier).
\end{itemize}

\subsubsection{2. Le Fournisseur}
Le fournisseur intervient comme cocontractant de l'administration. Ses interventions couvrent :
\begin{itemize}
  \item la participation à la mise en concurrence (consultation des expressions de besoin ouvertes à la soumission et dépôt d'une soumission) ;
  \item l'exécution matérielle de la commande (réalisation des livraisons et réception du procès-verbal établi par le comptable matière) ;
  \item la demande de règlement financier (introduction d'une demande de paiement accompagnée de la soumission des factures et des pièces justificatives) ;
  \item la consultation de l'avancement du traitement (suivi de l'évolution de ses demandes).
\end{itemize}

\subsubsection{3. Le Contrôleur financier}
Le contrôleur financier garantit la régularité budgétaire et financière des opérations à deux moments clés du processus :
\begin{itemize}
  \item \textbf{lors de l'engagement juridique :} il effectue la réservation des crédits, donne le quitus sur le classement des matières dans les natures économiques, et examine puis vise les projets d'actes juridiques soumis par l'ordonnateur ;
  \item \textbf{lors de l'ordonnancement :} il reçoit les projets d'ordonnancement signés par l'ordonnateur, effectue les contrôles nécessaires, puis signe et transmet les dossiers d'ordonnancement au comptable public.
\end{itemize}

\subsubsection{4. Le Comptable matière}
Le comptable matière assure le contrôle physique et patrimonial des biens livrés :
\begin{itemize}
  \item il prend en charge les livraisons effectuées par le fournisseur ;
  \item il constate les produits effectivement livrés ;
  \item il informe l'ordonnateur des livraisons réalisées ;
  \item il établit et délivre au fournisseur un procès-verbal de réception.
\end{itemize}

\subsubsection{5. Le Comptable public}
Dans les limites des informations communiquées à ce stade :
\begin{itemize}
  \item il reçoit les dossiers d'ordonnancement signés et transmis par le contrôleur financier.
\end{itemize}
\emph{Nota : conformément aux règles méthodologiques retenues, aucune responsabilité supplémentaire (telle que le contrôle de régularité comptable ou le décaissement) n'est ajoutée sans confirmation préalable du responsable du projet.}

\subsection{Tableau synthétique des acteurs et de leurs responsabilités}

Le tableau \ref{tab:acteurs_synthese} présente la synthèse structurée des cinq acteurs, en distinguant leur rôle métier global, leurs responsabilités métier, les opérations qu'ils réalisent dans le système et les acteurs avec lesquels ils interagissent.

\begin{landscape}
\small
\begin{longtable}{L{2.8cm} L{3.6cm} L{5.8cm} L{6.2cm} L{3.8cm}}
\caption{Tableau synthétique des acteurs, rôles, responsabilités métier, opérations et interactions}
\label{tab:acteurs_synthese} \\
\toprule
\textbf{\sffamily Acteur} &
\textbf{\sffamily Rôle dans le processus} &
\textbf{\sffamily Principales responsabilités métier} &
\textbf{\sffamily Principales opérations réalisées} &
\textbf{\sffamily Autres acteurs en interaction} \\
\midrule
\endfirsthead

\multicolumn{5}{l}{\footnotesize\itshape Suite du tableau \ref{tab:acteurs_synthese} — Acteurs et responsabilités} \\
\toprule
\textbf{\sffamily Acteur} &
\textbf{\sffamily Rôle dans le processus} &
\textbf{\sffamily Principales responsabilités métier} &
\textbf{\sffamily Principales opérations réalisées} &
\textbf{\sffamily Autres acteurs en interaction} \\
\midrule
\endhead

\bottomrule
\endfoot

\textbf{Ordonnateur} &
Initiateur de la dépense, autorité d'engagement juridique, de liquidation et d'ordonnancement &
\begin{itemize}[leftmargin=1em, nosep]
  \item Définition officielle du besoin
  \item Engagement juridique de l'organisme après visa préalable
  \item Vérification de la réalité de la dette et arrêté du montant exact de la dépense (liquidation)
  \item Émission et signature de l'ordre de payer (projet d'ordonnancement)
\end{itemize} &
\begin{itemize}[leftmargin=1em, nosep]
  \item Exprimer un besoin
  \item Préparer le projet d'acte juridique
  \item Soumettre le projet au contrôleur financier
  \item Signer l'acte juridique après les contrôles requis
  \item Recevoir les informations relatives aux livraisons
  \item Enregistrer et valider les factures
  \item Procéder à la liquidation
  \item Générer un projet d'ordonnancement de paiement
  \item Signer et transmettre le projet au contrôleur financier
\end{itemize} &
\begin{itemize}[leftmargin=1em, nosep]
  \item Fournisseur
  \item Contrôleur financier
  \item Comptable matière
\end{itemize} \\
\midrule

\textbf{Fournisseur} &
Opérateur économique soumissionnaire, exécutant des livraisons et créancier &
\begin{itemize}[leftmargin=1em, nosep]
  \item Formulation d'une offre commerciale en réponse à un besoin
  \item Exécution conforme des livraisons convenues
  \item Justification de sa créance en vue du paiement
\end{itemize} &
\begin{itemize}[leftmargin=1em, nosep]
  \item Consulter les expressions de besoin ouvertes à la soumission
  \item Soumissionner à une expression de besoin
  \item Effectuer les livraisons
  \item Introduire une demande de paiement
  \item Soumettre les factures et les pièces justificatives
  \item Suivre l'évolution de ses demandes
\end{itemize} &
\begin{itemize}[leftmargin=1em, nosep]
  \item Ordonnateur
  \item Comptable matière
\end{itemize} \\
\midrule

\textbf{Contrôleur financier} &
Autorité de contrôle budgétaire a priori sur l'engagement et l'ordonnancement &
\begin{itemize}[leftmargin=1em, nosep]
  \item Garantie de la disponibilité et du blocage des crédits budgétaires
  \item Validation de l'exactitude de l'imputation en nature économique (quitus)
  \item Contrôle de régularité et visa préalable des projets d'actes juridiques
  \item Contrôle final, signature et transmission des dossiers d'ordonnancement
\end{itemize} &
\begin{itemize}[leftmargin=1em, nosep]
  \item Effectuer la réservation des crédits
  \item Donner le quitus sur le classement des matières dans les natures économiques
  \item Examiner et viser les projets d'actes juridiques
  \item Recevoir les projets d'ordonnancement
  \item Effectuer les contrôles nécessaires
  \item Signer et transmettre les dossiers au comptable public
\end{itemize} &
\begin{itemize}[leftmargin=1em, nosep]
  \item Ordonnateur
  \item Comptable public
\end{itemize} \\
\midrule

\textbf{Comptable matière} &
Garant de la réception physique et de la constatation des biens livrés &
\begin{itemize}[leftmargin=1em, nosep]
  \item Réception matérielle et vérification effective des biens livrés
  \item Attestation officielle de la livraison par procès-verbal
  \item Information de l'ordonnateur pour permettre la liquidation
\end{itemize} &
\begin{itemize}[leftmargin=1em, nosep]
  \item Prendre en charge les livraisons
  \item Constater les produits effectivement livrés
  \item Informer l'ordonnateur des livraisons
  \item Établir et délivrer au fournisseur un procès-verbal de réception
\end{itemize} &
\begin{itemize}[leftmargin=1em, nosep]
  \item Fournisseur
  \item Ordonnateur
\end{itemize} \\
\midrule

\textbf{Comptable public} &
Destinataire institutionnel des dossiers d'ordonnancement validés &
\begin{itemize}[leftmargin=1em, nosep]
  \item Réception officielle des dossiers d'ordonnancement transmis après contrôle financier (dans la limite du périmètre documenté)
\end{itemize} &
\begin{itemize}[leftmargin=1em, nosep]
  \item Recevoir les dossiers d'ordonnancement transmis par le contrôleur financier
\end{itemize} &
\begin{itemize}[leftmargin=1em, nosep]
  \item Contrôleur financier
\end{itemize} \\
\end{longtable}
\end{landscape}
